\section{Experiments}
\label{sec:experiments}

\subsection{Setup}

\paragraph{Data and benchmark.} We evaluate on a dual-market benchmark. The US market pool consists of 25 representative Dow-30 equities spanning technology, consumer, energy, financial, and healthcare sectors, with OHLCV from yfinance, news from Alpaca and Reuters, and filings from SEC EDGAR. The China market supplement draws 10--15 high-liquidity CSI-300 constituents from Tushare. Training and self-evolution rollouts use 2019-01 to 2022-12; validation uses 2023-01 to 2023-12; the primary test interval is 2025-01 to 2025-12; and a live-forward window of 2026-05 to 2026-08 is reserved for post-submission evidence.

\paragraph{Baselines.} We compare against twelve baselines in four categories. \emph{Numerical time-series models:} PatchTST, TimesNet, iTransformer. \emph{LLM financial multi-agent systems:} TradingAgents, FinCon, FinAgent, FinRobot, R\&D-Agent-Quant, AlphaAgent. \emph{Factor mining baselines:} CogAlpha, FactorMAD. \emph{Market benchmarks:} Buy \& Hold, equal-weight portfolio, 60-40, SMA(50,200) crossover, pure momentum.

\paragraph{Metrics.} We report metrics at three levels. \emph{Prediction:} directional accuracy, IC, RankIC. \emph{Portfolio:} Sharpe ratio (with 15bps round-trip cost, 5bps slippage, 1-day execution delay), annualized return, maximum drawdown, Calmar, Sortino, win rate, profit factor. \emph{Self-evolution:} evolution curve (Sharpe versus iteration $k$), belief store size and hit rate, alpha decay versus $k$. All metrics are reported conditioned on regime (calm/normal/volatile $\times$ trending/ranging).

\subsection{Main Results}

[To be filled with experimental results upon completion of Stage 5.]

\subsection{Ablation Studies}

We conduct ten ablation experiments (A1--A10) to validate each component:

\begin{table}[h]
\centering
\small
\begin{tabular}{llp{6cm}}
\toprule
ID & Configuration & Hypothesis Tested \\
\midrule
A1 & Full \textsc{FinOPD} & Upper bound \\
A2 & Remove visual-geometric encoder & H1: value of geometric perception \\
A3 & Freeze VLM (no LoRA) & Value of financial chart alignment \\
A4 & Remove router, use fixed tool templates & H2: value of learnable routing \\
A5 & Disable parametric self-evolution (verbal reinforcement only) & H3: value of OPSD \\
A6a & Disable Shapley credit (uniform weights) & Value of credit assignment \\
A6b & Disable non-parametric belief consolidation & H3: value of long-term memory \\
A7 & Disable alignment guard & Necessity of drift protection \\
A8 & Single-agent collapse (5 $\to$ 1) & Value of multi-agent decomposition \\
A9 & Single iteration ($k{=}1$) vs.\ multi-iteration ($k{\geq}3$) & H4: evolution monotonicity \\
A10 & Replace OPSD with trajectory-search self-evolution & Method selection justification \\
\bottomrule
\end{tabular}
\caption{Ablation configurations and corresponding hypotheses.}
\label{tab:ablations}
\end{table}

\subsection{Evolution Curve Analysis}

[To be filled with evolution curve figures showing Sharpe vs.\ iteration $k$ across 3 seeds, demonstrating monotonic improvement under post-cutoff evaluation.]

\subsection{Evaluation Rigor}

We enforce three evaluation disciplines. \emph{Cutoff-aware rollout:} all test dates are strictly later than each model's pretraining cutoff. \emph{Counterfactual perturbation:} we apply three perturbation types (news sentiment polarity inversion, non-critical financial figure randomization, date token replacement) and report the Sharpe/accuracy delta $\Delta$; small $\Delta$ would indicate reliance on memorization rather than reasoning. \emph{Live-forward track:} a daemon pulls real-time data daily for inference without retraining, accumulating evidence for the rebuttal period.
